\documentclass[12pt]{article}

% graphicx package, useful for including eps and pdf graphics
\usepackage{graphicx}
\DeclareGraphicsExtensions{.png,.png,.jpg}

% basic packages
\usepackage{fancyhdr}
\usepackage{color}
\usepackage{parskip}
\usepackage{float}
\usepackage{microtype}
\usepackage{url}
\usepackage{hyperref}
\usepackage{booktabs}
\usepackage{makecell}
\usepackage{multirow}

% text layout
\usepackage[top=1in, bottom=1in, left=1.25in, right=1.25in, includefoot, heightrounded]{geometry}
\geometry{textwidth=17cm} % 15.25cm for single-space, 16.25cm for double-space
\geometry{textheight=22.5cm} % 22cm for single-space, 22.5cm for double-space

\newcommand{\PageLine}{\rule{\textwidth}{0.25mm}}
\newcommand{\VeryLongSignatureLine}{\rule{12cm}{0.25mm}}
\newcommand{\LongSignatureLine}{\rule{9cm}{0.25mm}}
\newcommand{\ShortSignatureLine}{\rule{7cm}{0.25mm}}
\newcommand{\DateLine}{\rule{3cm}{0.25mm}}

% helps to keep figures from being orphaned on a page by themselves
\renewcommand{\topfraction}{0.85}
\renewcommand{\textfraction}{0.1}

% bold the 'Figure #' in the caption and separate it with a period
% Captions will be left justified
\usepackage[labelfont=bf,labelsep=period,font=small]{caption}

% cite package, to clean up citations in the main text. Do not remove.
\usepackage{cite}

\setlength{\headheight}{67pt}
\pagestyle{fancy}
\fancyhf{}

\lhead{\includegraphics{logo}}
\rhead{\includegraphics{contact}}

\begin{document}

\textbf{NWGC SARS-CoV-2 (COVID-19) Qual PCR Lab Report}

\textbf{Participant Name: \LongSignatureLine}

\textbf{Specimen identifier: \LongSignatureLine}

\textbf{Date Sample Submitted: \ShortSignatureLine}

\textbf{Date Sample Tested: \ShortSignatureLine}

\textbf{Date results provided: Date \DateLine}

\textbf{Specimen Type:} Midnasal, self-collected

\textbf{\VeryLongSignatureLine}

\textbf{Result of SARS-CoV-2 (COVID-19) Qual PCR}

(ONLY ONE PROVIDED: Not detected / SARS-CoV-2 Detected (positive) / Inconclusive / Unable to be tested

\PageLine

\textbf{Please Note:}

\begin{itemize}
\item

  \textbf{The type of sample and testing used in this project is
  conducted as part of a public health surveillance program in
  partnership with King County and other agencies.} The testing you have
  received is not a substitute for receiving personal medical care by a
  healthcare provider.

\item

  \textbf{Specimen Chain of Custody:} the specimen tested was
  self-collected, self-labeled, and submitted to SCAN by mail or
  delivery service. It was not collected, labeled, and transported by a
  licensed healthcare professional. This may impact test results and
  false negative results cannot be excluded. Follow-up clinical testing
  may be warranted.

\item

  \textbf{Specimen collection device:} the mid turbinate swab is not an
  FDA/CDC Emergency Use Authorization (EUA) approved specimen type.
  Therefore, not detected/negative results may not exclude SARS-CoV-2
  infection.

\item

  \textbf{``Positive'' results may be verified by the Washington State
  Department of Health}. This result should be discussed with your
  healthcare provider if you have questions.

\item{
  \textbf{If your test says ``Not detected'' it means that SARS-CoV-2
    was not detected in our laboratory.}

  No test is perfect and there is a small chance that you do have the
  SARS-CoV-2 virus that causes COVID-19 with a negative test. This test
  does not use the standard way of collecting specimens by a health care
  provider, and the sample was sent by mail, not collected and transported
  in a healthcare setting. Regardless of your test result, it's best to
  stay home when you feel sick, and avoid contact with other people. If
  you have other health problems or are very sick or are worried about
  your illness, you should contact your healthcare provider. If your
  symptoms persist, you should discuss getting another test for COVID-19
  with your health care provider.
}

\item

  \textbf{If your test says ``Inconclusive''} it means that our
  laboratory test results for the SARS-CoV-2 target are not certain.

  No test is perfect and there is a chance that an inconclusive result
  could mean that you do have the SARS-CoV-2 virus that causes COVID-19.
  This test does not use the standard way of collecting specimens by a
  health care provider, and the sample was sent by mail, not collected and
  transported in a healthcare setting. Regardless of your test result,
  it's best to stay home when you feel sick, and avoid contact with other
  people. If you have other health problems or are very sick or are
  worried about your illness, you should contact your healthcare provider.
  If your symptoms persist, you should discuss getting another test for
  COVID-19 with your health care provider.

\item

  \textbf{If your test says ``Unable to be tested''} it means that there
  was either not enough of the specimen or there was an error in the
  test process.

  Since testing was unable to be performed, there is a chance that you do
  have the SARS-CoV-2 virus that causes COVID-19. This test does not use
  the standard way of collecting specimens by a health care provider, and
  the sample was sent by mail, not collected and transported in a
  healthcare setting. Regardless of your test result, it's best to stay
  home when you feel sick, and avoid contact with other people. If you
  have other health problems or are very sick or are worried about your
  illness, you should contact your healthcare provider. If your symptoms
  persist, you should discuss getting another test for COVID-19 with your
  health care provider.

\end{itemize}

\begin{itemize}
\item

  \textbf{Positive tests (shown as ``DETECTED'' or ``Virus Detected'')
  will be reported to the appropriate state and county Health Department
  by the testing facility.} You DO NOT need to report these results
  yourself. Public Health may follow up with you if your test is
  positive for novel coronavirus to ask you about your illness and
  people you may have come in contact with.

\item

  \textbf{You can print out or email these test results to your health
  care provider}

\end{itemize}

\textbf{Methodology of SARS-CoV-2 (COVID-19) Qual PCR:}

This assay is designed to detect SARS-CoV-2 using reverse-transcription
PCR. This assay relies on amplification of the viral genome segments
Orf1B and S with an internal human RNaseP amplification control. This
test was developed and its performance characteristics determined by the
University of Washington's Northwest Genomics Center. This laboratory
developed test is approved by Washington Department of Health under
emergency use authorization, but has not been approved by FDA. This
laboratory is certified under the Clinical Laboratory Improvement
Amendments (CLIA) as qualified to perform high complexity clinical
laboratory testing.

For further provider information, please see these links provided by the
Washington State Department of Health:

\begin{itemize}
\item

  \href{https://www.doh.wa.gov/Portals/1/Documents/1600/coronavirus/COVIDcasepositive.pdf}{\textbf{{Patients
  with confirmed or suspected COVID-19}}}

\item

  \href{https://www.doh.wa.gov/Portals/1/Documents/1600/coronavirus/COVIDexposed.pdf}{\textbf{{Patients
  who were exposed to a confirmed COVID-19 case}}}

\item

  \href{https://www.doh.wa.gov/Portals/1/Documents/1600/coronavirus/COVIDconcerned.pdf}{\textbf{{Unexposed
  patients with COVID-19 symptoms}}}

\end{itemize}

\textbf{Additional resources are available at the} Washington State
Department of Health \textbf{and CDC websites: }

\begin{itemize}
\item

  \href{https://www.doh.wa.gov}{\textbf{{https://www.doh.wa.gov/}}}

\item

  \href{https://www.cdc.gov/coronavirus/2019-ncov/index.html}{\textbf{{https://www.cdc.gov/coronavirus/2019-ncov/index.html}}}

\end{itemize}

\end{document}
